\section{Metrics and Validation}

The validation design asks several different questions. The realized path asks what
happened in the held-out monthly window. The point-in-time audit asks whether decisions
used only information available at the rebalance date. The regression asks whether the
option book is mostly equity drift or volatility exposure. The attribution asks whether P\&L
comes from delta, gamma, vega, carry, VIX-forward exposure, skew/tail terms, or residual
marking error. Cost stress asks whether the gross result survives realistic research
frictions.

\begin{figure}[H]
\centering
\IfFileExists{figures/portfolio_growth.pdf}
  {\includegraphics[width=0.88\textwidth]{portfolio_growth.pdf}}
  {\fbox{\parbox{0.82\textwidth}{Run the empirical pipeline to generate the portfolio growth figure.}}}
\caption{Headline out-of-sample growth paths. The figure focuses on the central option
allocators and equity benchmarks; the all-strategy growth diagnostic appears in
Appendix~A. The log-scale floor prevents failed strategies from disappearing visually.}
\label{fig:growth}
\end{figure}

Performance is reported with complementary measures. Sharpe follows the reward-to-
variability tradition \citep{sharpe1966}. Sortino replaces total volatility with downside
deviation \citep{sortino_price1994}. Calmar scales return by maximum drawdown
\citep{magdon_ismail_atiya2004}. Omega compares gains above a threshold with shortfalls
below it \citep{keating_shadwick2002}. Information ratio measures active return against
SPY \citep{grinold_kahn2000}. The zero monthly threshold is used for Sortino and Omega
because cash is a numeraire in the model, not an alpha sleeve.

The inference layer is intentionally conservative. Sharpe, Sortino, Calmar, Omega, regime
statistics, and leave-one-out differences are reported with block-bootstrap intervals where
meaningful. Factor alpha uses HAC/Newey--West standard errors. These intervals do not
prove economic significance, but they prevent the paper from reading point estimates as
certainties. The pre-production extension also reports probabilistic Sharpe ratios,
deflated Sharpe ratios, and a block-bootstrap max-Sharpe screen across the tested gross and
executable-cost variants.

\begin{figure}[H]
\centering
\IfFileExists{figures/random_sharpe_histogram.pdf}
  {\includegraphics[width=0.70\textwidth]{random_sharpe_histogram.pdf}}
  {\fbox{\parbox{0.72\textwidth}{Run the empirical pipeline to generate the random feasible Sharpe figure.}}}
\caption{Random feasible comparison. The histogram shows out-of-sample Sharpe ratios for
random gross-NAV-feasible option books; the vertical marker shows the optimized
Greek-Markowitz-plus-VIX sleeve. The exhibit asks whether optimization adds information
beyond arbitrary feasible option exposure.}
\label{fig:random_sharpe}
\end{figure}

The point-in-time contract is strict. Decision dates precede return dates. State snapshots,
Greeks, forecasts, covariances, residuals, betas, borrow proxies, and cost inputs are drawn
from the training window or the decision date. Expiry spot, future IV, future VIX/VVIX,
future VX curves, and future option returns are not allowed to enter portfolio construction.
Cost-scenario ledgers, required-capital denominators, hurdle/no-trade records, liquidity-tier
diagnostics, assignment-risk filters, and reality-check inputs obey the same rule. The
audit layer reads the machine-readable outputs independently and fails if core lineage,
settlement, cost, inference, claim, or manuscript checks fail.
